\chapter{Канонические барионные веса и запрет блочной перебалансировки}
\label{ch:v8-baryon-canonical-weights-geta-no-go-gate}

После закрытия временной ветви проверим первый переносимый барионный
результат. Рассматриваются не физические массы барионов, а две точные
алгебраические заявки: совпадение шести весов барионного расчёта с
каноническим представителем общего КМС-следа Тома VIII и недостижимость
внешней мишени вдоль положительного блочного семейства $G_\eta$.

\section{Точная точка общего КМС-следа}

Положим $x=e^{-2}$ и
\begin{equation}
 a=\frac1{11+10x},
 \qquad
 b=\frac{x}{11+10x},
 \qquad
 11a+10b=1,
 \qquad
 \frac ba=x.
 \label{eq:v8-baryon-kms-density}
\end{equation}
Следовые квадраты шести семейств дают веса
\begin{align}
 \kappa_L&=\frac1{13(a+b)}
 =\frac{10x+11}{13(x+1)},
 &
 \kappa_{3}&=\frac1{a+b}
 =\frac{10x+11}{x+1},
 \nonumber\\
 \kappa_{2}&=\frac1{(5a+b)/2}
 =\frac{20x+22}{x+5},
 &
 \kappa_{1}&=\frac1{(13a+25b)/6}
 =\frac{60x+66}{25x+13},
 \nonumber\\
 \kappa_Q&=\kappa_X=\frac1{2(a+b)}
 =\frac{10x+11}{2(x+1)}.
 \label{eq:v8-baryon-six-exact-weights}
\end{align}
Эти выражения покомпонентно совпадают с шестью коэффициентами
канонического шумового кадра общего КМС-следа. Максимальное расхождение с
сохранённой десятичной распечаткой равно $1.1\cdot10^{-14}$ и обусловлено
только преобразованием точного выражения к машинному числу.

\begin{proposition}[Провенанс шести весов]
Шесть весов барионного пакета принадлежат той же стабилизаторной точке
шестимерного конуса, которая была построена в Томе VIII ортонормировкой
шумового кадра относительно состояния
$\rho_N=aI_{11}\oplus bI_{10}$. Это тождество в $\mathbb Q(x)$, а не
рациональная аппроксимация десятичных коэффициентов.
\end{proposition}

Утверждение относится к стабилизаторной ветви при фиксированном фоне.
Оно не доказывает каноничность переноса направленных пар $V,V^*$ в
трёхкварковый расчёт и не повышает барионный дискриминатор до внутреннего
наблюдаемого Тома VIII.

\section{Блочное семейство G-эта}

В рассматриваемом барионном отображении общий множитель $\eta>0$ умножает переносные
семейства, оставляя калибровочные семейства неизменными. Дискриминатор
имеет точную форму
\begin{equation}
 v_\eta(x,\eta)=
 \frac{104(25x^2+38x+13)}
 {425\eta x^3+2346\eta x^2+1105\eta x
  +3692x^2+12584x+5564}.
 \label{eq:v8-baryon-geta-value}
\end{equation}
Его производная равна
\begin{equation}
 \frac{\partial v_\eta}{\partial\eta}
 =-\frac{1768x(x+1)(x+5)(25x+13)^2}
 {\left(425\eta x^3+2346\eta x^2+1105\eta x
 +3692x^2+12584x+5564\right)^2}<0
 \label{eq:v8-baryon-geta-derivative}
\end{equation}
для всех $x>0$ и $\eta>0$. Поэтому максимум на положительном срезе есть
предел при $\eta\to0+$.

Внешняя сравнительная мишень записывается точно как
\begin{equation}
 v_{\rm ext}=\frac{293.1}{3\cdot349.0}=\frac{977}{3490}.
 \label{eq:v8-baryon-external-target}
\end{equation}
Решение уравнения $v_\eta=v_{\rm ext}$ имеет вид
\begin{equation}
 \eta_*(x)=
 \frac{52\left(105133x^2+28806x-13799\right)}
 {16609x(x+5)(25x+13)}.
 \label{eq:v8-baryon-geta-root}
\end{equation}

\section{Точный знак без десятичной подстановки}

Обозначим
\begin{equation}
 P(x)=105133x^2+28806x-13799.
 \label{eq:v8-baryon-sign-polynomial}
\end{equation}
На положительной полуоси $P'(x)=210266x+28806>0$. Кроме того,
\begin{equation}
 e^2>\sum_{n=0}^{4}\frac{2^n}{n!}=7,
 \qquad
 x=e^{-2}<\frac17,
 \qquad
 P\!\left(\frac17\right)=-\frac{52768}{7}<0.
 \label{eq:v8-baryon-exp-exact-bound}
\end{equation}
Следовательно, $P(e^{-2})<0$ и $\eta_*(e^{-2})<0$ строго. Эквивалентно,
\begin{equation}
 v_{\rm ext}-\sup_{\eta>0}v_\eta
 =-\frac{P(x)}{3490(71x^2+242x+107)}>0
 \quad\text{при }x=e^{-2}.
 \label{eq:v8-baryon-geta-strict-gap}
\end{equation}
Для чтения, но не для доказательства,
\begin{align}
 v_1&=0.256960942726249\ldots,\nonumber\\
 \sup_{\eta>0}v_\eta&=0.263742325124295\ldots,\nonumber\\
 \eta_*&=-2.19282832819175\ldots.
 \label{eq:v8-baryon-geta-display}
\end{align}

\begin{theorem}[Запрет положительной блочной перебалансировки]
В точном рассматриваемом барионном отображении при $x=e^{-2}$ внешняя мишень
$977/3490$ недостижима для всех $\eta>0$ вдоль семейства
$G_\eta=\eta P_{\rm tr}+P_{\rm gauge}$. Требуемый корень отрицателен, а
дискриминатор строго убывает по $\eta$.
\end{theorem}

\begin{nogocriterion}[Граница физического чтения]
Теорема запрещает только один точно заданный способ общей
перебалансировки переносных и калибровочных семейств. Внешняя мишень
содержит подгоночный массовый параметр, а перенос направленных компонент
шумового кадра в барионную форму ещё неоднозначен. Поэтому число
$8.21\%$ нельзя объявлять безусловным предсказанием Тома VIII или
физической теоремой о барионных массах.
\end{nogocriterion}

\section{Следующий различающий вопрос}

Два исходных чтения давали значения дискриминатора около $0.2570$ и
$0.1962$. Следующая глава восстанавливает из 19-мерного факторпространства общего КМС-следа
точное правило, по которому пары $V,V^*$ переходят в одночастичную и
парную трёхкварковые формы. До такого селектора дальнейшая подгонка
семейств скоростей не имеет однозначного объекта.

\begin{protocol}
\begin{enumerate}
 \item поле коэффициентов: \textbf{$\mathbb Q(x)$, $x=e^{-2}$};
 \item нормировка $11a+10b=1$: \textbf{точно};
 \item отношение $b/a=x$: \textbf{точно};
 \item шесть следовых весов: \textbf{воспроизведены точно};
 \item ветвь каноничности: \textbf{стабилизаторная};
 \item монотонность $v_\eta$: \textbf{строго убывает};
 \item доказательство $e^{-2}<1/7$: \textbf{точное};
 \item корень $\eta_*$: \textbf{отрицателен};
 \item мишень на всём $\eta>0$: \textbf{недостижима};
 \item внешний остаток $8.21\%$: \textbf{сравнение, не вывод};
 \item единственность направленного одночастичного ограничения:
 \textbf{доказана в следующей главе};
 \item физическая теорема о массах: \textbf{не получена};
 \item следующий гейт: \textbf{селектор направленного переноса, выполнен в следующей главе}.
\end{enumerate}
\end{protocol}

Машинный сертификат сохранён в
\path{s2t_v8_baryon_canonical_weights_geta_no_go_gate_results.json}.